\chapter{Przedziały ufności: konstrukcja, pokrycie i interpretacja}
\label{ch:confidence-intervals}
\paragraph{Powiązania i odsyłacze.} Związek przedziałów z testami wyjaśnia rozdział~\ref{ch:testing}; przykłady dokładnych rozkładów statystyk podaje rozdział~\ref{ch:parametric-tests}.


\section{Co dokładnie oznacza 95\% ufności?}
Zaobserwowany odsetek wadliwych produktów wynosi $7/100$. Samo
$\widehat p=0{,}07$ nie mówi, jaka jest precyzja wnioskowania.
Przedział ufności jest sposobem opisania tej niepewności, ale
nie jest rozkładem prawdopodobieństwa nieznanego parametru.

\begin{definicja}[Zbiór ufności i pokrycie]\label{def:confidence-set}
Niech $(\mathsf X,\mathcal A,\{P_\theta:\theta\in\Theta\})$
będzie eksperymentem statystycznym, a
$g:\Theta\to\mathsf T$ ustaloną funkcją parametru,
gdzie $(\mathsf T,\mathcal T)$ jest przestrzenią mierzalną.
\emph{Zbiorem ufności} na poziomie $1-\alpha$, gdzie
$0<\alpha<1$, nazywamy losowy zbiór $C(X)\subseteq\mathsf T$
taki, że dla każdego $\theta\in\Theta$ zdarzenie
$\{g(\theta)\in C(X)\}$ jest mierzalne i
\[
 P_\theta\bigl(g(\theta)\in C(X)\bigr)\ge1-\alpha.
\]
Jeśli $\mathsf T=\R$ i $C(X)=[L(X),U(X)]$ dla
mierzalnych $L\le U$, nazywamy go \emph{przedziałem ufności}.
Powyższe prawdopodobieństwo jest \emph{pokryciem} przy $\theta$.
Gdy dla każdego parametru zachodzi równość, mówimy o
\emph{dokładnym pokryciu nominalnym}; procedura z pokryciem
większym od nominalnego jest \emph{konserwatywna}.
\end{definicja}

Losowość należy do przedziału $C(X)$, zanim zaobserwujemy próbę.
Po podstawieniu konkretnych danych $x$ otrzymujemy ustalony zbiór
$C(x)$. W interpretacji klasycznej nie stwierdzamy wtedy, że
„parametr z prawdopodobieństwem $95\%$ leży w tym konkretnym
przedziale”. Stwierdzamy, że \emph{procedura} konstruowania
przedziałów ma zadane pokrycie w powtarzanych doświadczeniach.
Rozkład wiarygodności a posteriori prowadzący do innego rodzaju
przedziału zostanie wprowadzony w części bayesowskiej.

\begin{uwaga}[Pokrycie punktowe, jednostajne i przybliżone]
Zapis $P_\theta(g(\theta)\in C_n)\to1-\alpha$ dla każdego
ustalonego $\theta$ oznacza pokrycie asymptotyczne \emph{punktowe}.
Silniejsze żądanie
$\liminf_{n\to\infty}\inf_{\theta\in\Theta}
P_\theta(g(\theta)\in C_n)\ge1-\alpha$
gwarantuje pokrycie asymptotyczne \emph{jednostajne}.
Pierwsze nie implikuje drugiego; przy granicach przestrzeni
parametrów klasyczne przybliżenie normalne może działać
bardzo źle. W zastosowaniach małych prób pytamy zatem nie tylko
o nominalny procent, lecz także o zakres parametrów, dla których
gwarancja jest prawdziwa.
\end{uwaga}

\section{Dualność testów i przedziałów ufności}
\begin{twierdzenie}[Odwracanie rodziny testów]\label{thm:test-inversion}
Niech $\{P_\theta:\theta\in\Theta\}$ będzie rodziną miar
na $(\mathsf X,\mathcal A)$ oraz $g:\Theta\to\mathsf T$.
Dla każdego $t\in g(\Theta)$ niech $\varphi_t:\mathsf X\to\{0,1\}$
będzie mierzalnym testem hipotezy
$H_{0,t}:g(\theta)=t$ o poziomie co najwyżej $\alpha$:
\[
 \sup_{\theta:g(\theta)=t}\E_\theta\varphi_t(X)\le\alpha.
\]
Zakładamy ponadto, że zbiór
$C(x)=\{t\in g(\Theta):\varphi_t(x)=0\}$ ma mierzalne
zdarzenia przynależności. Wtedy $C(X)$ jest zbiorem ufności
dla $g(\theta)$ o pokryciu co najmniej $1-\alpha$.
\end{twierdzenie}
\begin{proof}
Ustalmy $\theta\in\Theta$ i przyjmijmy $t=g(\theta)$.
Z definicji konstrukcji $g(\theta)\notin C(X)$ wtedy i tylko
wtedy, gdy $\varphi_t(X)=1$. Zatem
$P_\theta(g(\theta)\notin C(X))
=\E_\theta\varphi_t(X)\le\alpha$, co daje tezę.
Nie trzeba zakładać, że test dla każdego $t$ odrzuca
jednakowe wartości obserwowanej statystyki: istotna jest
\emph{jednoczesna definicja rodziny testów}, a nie pojedynczy
test dopasowany do danych.
\end{proof}

\begin{uwaga}[Co oznacza odwrotna implikacja?]
Dla danego zbioru ufności $C(X)$ można zdefiniować test
hipotezy $g(\theta)=t$ przez
$\varphi_t(X)=\ind{\{t\notin C(X)\}}$. Jeśli
pokrycie jest co najmniej $1-\alpha$, każdy z tych testów ma
poziom co najwyżej $\alpha$. Testy i przedziały są więc
różnymi reprezentacjami tej samej reguły wnioskowania,
\emph{o ile} zostały skonstruowane z tą samą rodziną hipotez
i tymi samymi założeniami.
\end{uwaga}

\section{Prawidłowy model normalny: przedział dla średniej i wariancji}
\begin{twierdzenie}[Średnia i wariancja przy nieznanych parametrach]
\label{thm:normal-mean-var-ci}
Niech $n\ge2$, $X_1,\ldots,X_n\iid\Normal(\mu,\sigma^2)$,
$\mu\in\R$, $\sigma>0$, oraz niech $\bar X$ i $S^2$ oznaczają
średnią i wariancję z mianownikiem $n-1$. Dla $0<\alpha<1$
przedziały
\[
 C_\mu(X)=\left[\bar X-t_{n-1,1-\alpha/2}\frac{S}{\sqrt n},
 \bar X+t_{n-1,1-\alpha/2}\frac{S}{\sqrt n}\right],
\]
\[
 C_{\sigma^2}(X)=\left[
 \frac{(n-1)S^2}{\chi^2_{n-1,1-\alpha/2}},
 \frac{(n-1)S^2}{\chi^2_{n-1,\alpha/2}}\right]
\]
mają \emph{każdy osobno} pokrycie dokładnie $1-\alpha$.
Symbol $\chi^2_{\nu,q}$ oznacza kwantyl rzędu $q$ rozkładu
chi-kwadrat o $\nu>0$ stopniach swobody; symbol
$t_{\nu,q}$ wprowadzono w rozdziale~\ref{ch:parametric-tests}.
\end{twierdzenie}
\begin{proof}
Rozkład próby normalnej zapewnia niezależne zmienne
$Z=\sqrt n(\bar X-\mu)/\sigma\sim\Normal(0,1)$ i
$V=(n-1)S^2/\sigma^2\sim\chi^2_{n-1}$. Stąd
$Z/\sqrt{V/(n-1)}\sim t_{n-1}$. Zdarzenie
$\{|Z/\sqrt{V/(n-1)}|\le t_{n-1,1-\alpha/2}\}$ jest
dokładnie zdarzeniem $\{\mu\in C_\mu(X)\}$ i ma
prawdopodobieństwo $1-\alpha$. Analogicznie
$\{\chi^2_{n-1,\alpha/2}\le V\le
\chi^2_{n-1,1-\alpha/2}\}$ ma prawdopodobieństwo
$1-\alpha$ i po rozwiązaniu nierówności względem
$\sigma^2$ daje drugi przedział. Użyliśmy ciągłości
rozkładów; dla obu przedziałów razem nie zadeklarowano
tu pokrycia jednoczesnego.
\end{proof}

\begin{stwierdzenie}[Przedział dla ilorazu dwóch wariancji]
\label{prop:ratio-var-ci}
Niech próby $X_1,\ldots,X_n\iid\Normal(\mu_X,\sigma_X^2)$
i $Y_1,\ldots,Y_m\iid\Normal(\mu_Y,\sigma_Y^2)$ będą
niezależne, $n,m\ge2$, $\sigma_X,\sigma_Y>0$.
Jeśli $R=S_X^2/S_Y^2$ oraz $q=\sigma_X^2/\sigma_Y^2$,
to przedział
\[
 \left[\frac{R}{F_{n-1,m-1,1-\alpha/2}},
       \frac{R}{F_{n-1,m-1,\alpha/2}}\right]
\]
obejmuje $q$ z prawdopodobieństwem $1-\alpha$,
gdzie $F_{a,b,u}$ jest kwantylem rzędu $u$
rozkładu Fishera--Snedecora z $a,b>0$ stopniami swobody.
\end{stwierdzenie}
\begin{proof}
Zmienne
$(n-1)S_X^2/\sigma_X^2$ i
$(m-1)S_Y^2/\sigma_Y^2$ są niezależne i mają
rozkłady chi-kwadrat. Z definicji rozkładu F
$R/q\sim F_{n-1,m-1}$, więc odwrócenie nierówności
między kwantylami daje wskazany przedział.
\end{proof}

\begin{przyklad}[Dwie niezależne technologie a pomiary sparowane]
W zadaniu~\ref{zad:welch-case} dane pochodzą z dwóch
\emph{niezależnych} grup. Przy nierównych wariancjach
przedział dla $\mu_X-\mu_Y$ uzyskany metodą Welcha ma
postać
$[\bar X-\bar Y\pm t_{\nu,1-\alpha/2}
\sqrt{S_X^2/n+S_Y^2/m}]$, gdzie losowe $\nu$
zdefiniowano w rozdziale~\ref{ch:parametric-tests}.
Jest to przedział \emph{przybliżony}, nie dokładny.
Jeżeli te same jednostki zmierzono dwukrotnie i różnice
$D_i$ są niezależne między parami oraz normalne,
przedział oparty na $\bar D\pm t_{n-1,1-\alpha/2}S_D/\sqrt n$
jest \emph{dokładny}. Nie wolno podmieniać tych procedur
bez zmiany założeń o zależności.
\end{przyklad}

\section{Odsetek wad: trzy przedziały dla parametru Bernoulliego}
Niech $X_1,\ldots,X_n\iid\operatorname{Bernoulli}(p)$,
$n\ge1$, $p\in[0,1]$, $S=\sum_iX_i$ i
$\widehat p=S/n$. Trzy popularne konstrukcje odpowiadają
na to samo pytanie, ale opierają się na różnych argumentach.

\paragraph{Wald: prosta formuła bez gwarancji przy małych częstościach.}
Dla $z=z_{1-\alpha/2}$ przedział Walda ma postać
\[
 C_{\rm Wald}(X)=\left[\widehat p-z\sqrt{\frac{\widehat p(1-\widehat p)}n},
\widehat p+z\sqrt{\frac{\widehat p(1-\widehat p)}n}\right]
\cap[0,1].
\]
Z twierdzenia centralnego granicznego i twierdzenia Slutsky'ego
wynika pokrycie zbieżne do $1-\alpha$ \emph{dla każdego
ustalonego} $p\in(0,1)$. Nie jest to jednak gwarancja
jednostajna, szczególnie blisko $p=0$ i $p=1$.
Przy $S=0$ przedział redukuje się do $\{0\}$, więc dla
każdego prawdziwego $p>0$ ten wynik próby nie pokrywa
parametru. W praktyce nie zalecamy Walda dla rzadkich
zdarzeń ani małych prób.

\begin{twierdzenie}[Przedział Wilsona przez odwrócenie testu score]
\label{thm:wilson-ci}
Przy powyższych założeniach i $0<\alpha<1$ niech
$z=z_{1-\alpha/2}$. Zbiór parametrów $p\in[0,1]$
spełniających
\[
 \frac{n(\widehat p-p)^2}{p(1-p)}\le z^2,
 \qquad 0<p<1,
\]
z domknięciem rozwiązań przy $p=0,1$, jest przedziałem
$[c-h,c+h]$, gdzie
\[
 c=\frac{\widehat p+z^2/(2n)}{1+z^2/n},\qquad
 h=\frac{z}{1+z^2/n}\sqrt{
 \frac{\widehat p(1-\widehat p)}n+\frac{z^2}{4n^2}}.
\]
To \emph{przedział Wilsona}. Jego pokrycie wynosi
$1-\alpha+o(1)$ przy $n\to\infty$ dla każdego
ustalonego $p\in(0,1)$, lecz nie musi wynosić dokładnie
$1-\alpha$ dla skończonego $n$.
\end{twierdzenie}
\begin{proof}
W modelu Bernoulliego dla ustalonego $p\in(0,1)$
statystyka score $Z_p=\sqrt n(\widehat p-p)/\sqrt{p(1-p)}$
ma na mocy CLT rozkład asymptotycznie standardowy normalny.
Nierówność $|Z_p|\le z$ jest równoważna
$n(\widehat p-p)^2\le z^2p(1-p)$, czyli
\[
 (n+z^2)p^2-(2n\widehat p+z^2)p+n\widehat p^2\le0.
\]
Trójmian kwadratowy o dodatnim współczynniku przy $p^2$
przyjmuje wartości niedodatnie między pierwiastkami.
Wzory na środek i połowę ich różnicy upraszczają się
odpowiednio do $c$ i $h$. Dla $\widehat p=0$ albo $1$
odpowiednie domknięcie w punkcie brzegowym daje tę samą
formułę; otrzymane końce zawsze należą do $[0,1]$.
Przy ustalonym $p\in(0,1)$ na mocy CLT prawdopodobieństwo
nierówności dąży do $1-\alpha$. Nie wyciągamy wniosku
o jednostajnym pokryciu bez dodatkowego dowodu.
\end{proof}

\begin{twierdzenie}[Dokładny konserwatywny przedział Cloppera--Pearsona]
\label{thm:clopper-pearson}
Dla $0<\alpha<1$, $s\in\{0,\ldots,n\}$ definiujemy
$L(s)$ i $U(s)$ następująco:
\begin{itemize}
\item $L(0)=0$, a dla $s\ge1$ liczba $L(s)\in(0,1)$
jest jedynym rozwiązaniem $P_{L(s)}(S\ge s)=\alpha/2$;
\item $U(n)=1$, a dla $s\le n-1$ liczba $U(s)\in(0,1)$
jest jedynym rozwiązaniem $P_{U(s)}(S\le s)=\alpha/2$,
\end{itemize}
gdzie w obu równaniach $S\sim\Bin(n,p)$ z parametrem
$p$ wskazanym w indeksie. Wtedy dla każdego $p\in[0,1]$
\[
 P_p\bigl(L(S)\le p\le U(S)\bigr)\ge1-\alpha.
\]
\end{twierdzenie}
\begin{proof}
Dla $1\le s\le n$ funkcja
$p\mapsto P_p(S\ge s)$ jest ciągła i ściśle rosnąca
na $(0,1)$, z granicami $0$ i $1$; wynika to na przykład
z różniczkowania sumy dwumianowej albo ze sprzężenia
wskaźnikami $\ind{\{U_i\le p\}}$ z rozdziału~\ref{ch:testing}.
Analogicznie $p\mapsto P_p(S\le s)$ dla
$0\le s\le n-1$ jest ściśle malejąca od $1$ do $0$.
Zatem końce istnieją i są jednoznaczne.

Ustalmy prawdziwe $p\in(0,1)$ i oznaczmy przez
$S'$ pomocniczą zmienną o rozkładzie $\Bin(n,p)$.
Jeśli $p<L(S)$, to z monotoniczności ogona w parametrze
otrzymujemy $P_p(S'\ge S)<\alpha/2$.
Dla każdego dyskretnego rozkładu funkcja
$u(s)=P_p(S'\ge s)$ jest nierosnąca, a zbiór
$\{s:u(s)\le t\}$ jest końcowym przedziałem liczb
całkowitych (albo jest pusty). Jeżeli jego najmniejszy
element to $s_t$, mamy
$P_p(u(S)\le t)=P_p(S\ge s_t)=u(s_t)\le t$.
Stąd $P_p(p<L(S))\le\alpha/2$.
Analogicznie używając lewego ogona $P_p(S'\le S)$
otrzymujemy $P_p(p>U(S))\le\alpha/2$.
Nierówność sumy prawdopodobieństw daje
$P_p(p\notin[L(S),U(S)])\le\alpha$.
Dla $p=0$ lub $p=1$ właściwy koniec przedziału jest
równy parametrowi z prawdopodobieństwem jeden.
\end{proof}

\begin{uwaga}[Dokładny test, lecz niekoniecznie pokrycie równe nominalnemu]
Słowo „dokładny” odnosi się tu do gwarancji skończonopróbkowej
bez przybliżenia normalnego. Wskutek dyskretności rozkładu
pokrycie Cloppera--Pearsona jest często \emph{większe} niż
$1-\alpha$, a przedział może być szerszy niż Wilsona.
Nie należy mówić, że „Clopper--Pearson zawsze pokrywa dokładnie
$95\%$” w znaczeniu równości dla każdego $p$.
\end{uwaga}

\begin{zadanie}[7 wad na 100 sztuk: trzy różne odpowiedzi]
\label{zad:ci-seven}
Dla $n=100$, $s=7$ i $\alpha=0{,}05$ oblicz trzy
przedziały dla $p$ i zdecyduj, czy istnieją mocne podstawy,
by przyjąć $p<0{,}1$.
\par\emph{Rozwiązanie.} Punktowy estymator wynosi
$\widehat p=0{,}07$. Dla $z_{0{,}975}\approx1{,}96$
po podstawieniu do wzorów otrzymujemy
\[
 \begin{aligned}
 C_{\mathrm{Wald}}&\approx[0{,}01999;\ 0{,}12001],\\
 C_{\mathrm{Wilson}}&\approx[0{,}03432;\ 0{,}13750],\\
 C_{\mathrm{CP}}&\approx[0{,}02861;\ 0{,}13892].
 \end{aligned}
\]
Przedziały różnią się istotnie dolnym końcem, choć
wszystkie obejmują wartość $0{,}1$. Dla hipotezy
$H_0:p\ge0{,}1$ przeciw $p<0{,}1$ z rozdziału~\ref{ch:testing}
\emph{dokładna jednostronna górna granica} poziomu $95\%$
to $U_1(s)$ rozwiązujące
$P_{U_1(s)}(S\le s)=0{,}05$; przy $s=7$
wynosi $U_1(7)\approx0{,}12746$. Wobec tego
nie możemy na poziomie $5\%$ stwierdzić, że
$p<0{,}1$. Równoważne jest to temu, że jednostronny
test dokładny daje $p\text{-value}\approx0{,}20605$.
\emph{Dwustronny} przedział ufności poziomu $95\%$
nie jest automatycznie dualny do \emph{jednostronnego}
testu na poziomie $5\%$; ich zgodność tutaj nie
usprawiedliwia utożsamiania obu procedur.
\end{zadanie}

\begin{zadanie}[Przedział dla średniej różnicy w pomiarach sparowanych]
W zadaniu z rozdziału~\ref{ch:parametric-tests} osiem
niezależnych par miało różnice
$D=(3,2,4,3,1,2,5,2)$. Pod założeniem normalności
różnic przedział $95\%$ dla średniej różnicy $\delta$
ma postać $\bar D\pm t_{7,0{,}975}S_D/\sqrt8$.
Ponieważ $\bar D=2{,}75$, $S_D^2=23/14$ i
$t_{7,0{,}975}\approx2{,}365$, otrzymujemy
przedział w przybliżeniu $[1{,}68;\ 3{,}82]$.
Wnioskujemy o średniej \emph{różnicy wewnątrz jednostki},
a nie o różnicy średnich z dwóch niezależnych prób.
Interpretacja przyczynowa nadal wymaga odpowiedniego
projektu badania.
\end{zadanie}

\section{Proces Poissona: przedział dla intensywności zdarzeń}
Jeśli w czasie ekspozycji $T>0$ zaobserwowano $S=s$
zdarzeń, a model zakłada $S\sim\Pois(\lambda T)$,
to estymator $\widehat\lambda=s/T$ nie opisuje niepewności
przy małych $s$. W szczególności $s=0$ nie dowodzi,
że $\lambda=0$.

\begin{twierdzenie}[Dokładny przedział dla intensywności Poissona]
\label{thm:poisson-ci}
Niech $T>0$, $S\sim\Pois(\lambda T)$, $\lambda\ge0$,
oraz $0<\alpha<1$. Zdefiniujmy $a(0)=0$,
$a(s)=\chi^2_{2s,\alpha/2}/2$ dla $s\ge1$, a
$b(s)=\chi^2_{2(s+1),1-\alpha/2}/2$ dla $s\ge0$.
Wówczas
\[
 C_\lambda(S)=\left[\frac{a(S)}T,\frac{b(S)}T\right]
\]
ma pokrycie co najmniej $1-\alpha$ dla wszystkich
$\lambda\ge0$.
\end{twierdzenie}
\begin{proof}
Dla $s\ge1$ i $u>0$ związek pomiędzy rozkładami
Poissona i chi-kwadrat daje
\[
 P_{u}(S\ge s)=P(\chi^2_{2s}\le2u),\qquad
 P_{u}(S\le s)=P(\chi^2_{2(s+1)}>2u),
\]
gdzie w ogonach po lewej $S\sim\Pois(u)$.
Tożsamości wynikają z całkowania gęstości gamma albo
z pochodnej odpowiednich skończonych sum
$e^{-u}\sum_{j=0}^{s-1}u^j/j!$ i
$e^{-u}\sum_{j=0}^{s}u^j/j!$ z warunkiem brzegowym
przy $u=0$. Stąd $a(s)$ rozwiązuje
$P_{a(s)}(S\ge s)=\alpha/2$, a $b(s)$ rozwiązuje
$P_{b(s)}(S\le s)=\alpha/2$.
Tak jak w dowodzie twierdzenia~\ref{thm:clopper-pearson},
przekroczenie dolnego lub górnego końca oznacza wystąpienie
odpowiedniej jednostronnej wartości p mniejszej niż
$\alpha/2$. Nierówność sumy prawdopodobieństw kończy
dowód. Dla $\lambda=0$ mamy $S=0$ prawie na pewno
i przedział zawiera zero.
\end{proof}

\begin{przyklad}[Zero awarii nie oznacza zerowej intensywności]
Przy $s=0$ dolny koniec przedziału wynosi zero,
a górny wynosi $-\log(\alpha/2)/T$. Dla
$\alpha=0{,}05$ i $T=100$ godzin daje to około
$0{,}03689$ awarii na godzinę. Jeśli z góry
interesuje nas wyłącznie \emph{górna} granica, można
uzyskać krótszy, jednostronny przedział z końcem
$-\log(\alpha)/T\approx0{,}02996$; trzeba jednak
podać w raporcie, że chodzi o granicę jednostronną.
\end{przyklad}

\section{Gdy nie ufamy rozkładowi: przedział z nierówności Hoeffdinga}
\begin{twierdzenie}[Przedział rozkładowo niezależny dla średniej]
\label{thm:hoeffding-ci}
Niech $X_1,\ldots,X_n$ będą niezależnymi zmiennymi
losowymi o wspólnej wartości oczekiwanej $\mu$,
a przy tym $a\le X_i\le b$ prawie na pewno dla
pewnych znanych liczb $a<b$. Dla $0<\alpha<1$
przedział
\[
 \left[\bar X-(b-a)\sqrt{\frac{\log(2/\alpha)}{2n}},
       \bar X+(b-a)\sqrt{\frac{\log(2/\alpha)}{2n}}\right]
\]
ma pokrycie co najmniej $1-\alpha$.
\end{twierdzenie}
\begin{proof}
Nierówność Hoeffdinga dla niezależnych zmiennych
ograniczonych z tym samym przedziałem wartości zapewnia
$P(|\bar X-\mu|\ge\varepsilon)
\le2\exp[-2n\varepsilon^2/(b-a)^2]$.
Po podstawieniu
$\varepsilon=(b-a)\sqrt{\log(2/\alpha)/(2n)}$
prawa strona jest równa $\alpha$. Ponieważ zdarzenie
niepokrycia przedziałem implikuje
$|\bar X-\mu|\ge\varepsilon$, otrzymujemy tezę.
Nie wymagana jest normalność, ale \emph{znane ograniczenie}
$a\le X_i\le b$ jest istotnym założeniem.
\end{proof}

\begin{uwaga}[Praktyczne ograniczenia]
Przedział Hoeffdinga bywa szerszy od przedziału t lub Wilsona,
ponieważ kupujemy gwarancję przy słabszych założeniach
o rozkładzie. Nie działa automatycznie dla nieograniczonego
czasu awarii, wielokrotnych pomiarów tej samej osoby ani
danych z nieznanym mechanizmem selekcji.
Przedział $1-\alpha$ dla jednego ustalonego $n$ nie jest
automatycznie \emph{sekwencją ufności} zachowującą pokrycie
przy dowolnym, wybranym adaptacyjnie czasie zatrzymania.
\end{uwaga}

\section{Jak wybrać przedział i jak go raportować?}
\begin{center}\small
\begin{tabular}{>{\raggedright\arraybackslash}p{0.31\textwidth}>{\raggedright\arraybackslash}p{0.29\textwidth}>{\raggedright\arraybackslash}p{0.30\textwidth}}
\toprule
\textbf{Cel i dane} & \textbf{Konstrukcja} & \textbf{Co trzeba sprawdzić?}\\
\midrule
Średnia przy iid normalnych, $\sigma$ nieznana & Przedział t & Niezależność, normalność modelu; przy małych $n$ szczególnie istotne.\\
Średnia w dużej próbie o skończonej wariancji & Asymptotyczny przedział ze średniej i SE & Ciężkie ogony, zależność, wiarygodność oszacowania SE.\\
Średnia przy znanych granicach obserwacji & Hoeffding & Czy każde $X_i$ naprawdę mieści się w $[a,b]$?\\
Prawdopodobieństwo sukcesu w próbie Bernoulliego & Wilson & Przybliżone pokrycie; ostrożność przy skrajnych $p$.\\
Prawdopodobieństwo, mała próba lub skrajne częstości & Clopper--Pearson & Konserwatywność, możliwie duża szerokość.\\
Wariancja/iloraz wariancji próby normalnej & Chi-kwadrat/F & Wrażliwość na brak normalności.\\
Różnica średnich dwóch niezależnych grup & Welch (przybliżony) & Niezależność grup, ogony, próby niezrównoważone.\\
Intensywność zliczeń i znana ekspozycja & Dokładny Poisson & Model stałej intensywności; nadmierna dyspersja.\\
\bottomrule
\end{tabular}
\end{center}

\paragraph{Co koniecznie podajemy?}
Podajemy estymowany parametr i jego jednostkę, schemat
losowania, zastosowaną konstrukcję, poziom ufności,
liczebność lub ekspozycję, punktowe oszacowanie oraz
oba końce przedziału. W metodach asymptotycznych warto
podać przybliżenie i przeprowadzić analizę wrażliwości.
Jeżeli porównujemy wiele efektów naraz, przedziały
\emph{indywidualne} nie zapewniają automatycznie
jednoczesnego pokrycia całej rodziny.

\begin{lstlisting}[language=Python,caption={Porównanie przedziałów Walda, Wilsona i Cloppera--Pearsona oraz granica dla Poissona (SciPy).}]
from math import sqrt
from scipy.stats import beta, chi2, norm

n, s, alpha = 100, 7, 0.05
phat = s / n
z = norm.ppf(1 - alpha / 2)
wald = (max(0.0, phat - z * sqrt(phat * (1-phat)/n)),
        min(1.0, phat + z * sqrt(phat * (1-phat)/n)))
center = (phat + z*z/(2*n)) / (1 + z*z/n)
half = z/(1 + z*z/n) * sqrt(phat*(1-phat)/n + z*z/(4*n*n))
wilson = (center - half, center + half)
clopper = (0.0 if s == 0 else beta.ppf(alpha/2, s, n-s+1),
           1.0 if s == n else beta.ppf(1-alpha/2, s+1, n-s))
print('Wald:', wald)
print('Wilson:', wilson)
print('Clopper-Pearson:', clopper)

s_events, exposure = 0, 100.0
rate_upper_2s = chi2.ppf(1-alpha/2, 2*s_events+2) / (2*exposure)
print('Poisson upper (two-sided):', rate_upper_2s)
\end{lstlisting}

\begin{lstlisting}[language=Julia,caption={Przedział Wilsona i dokładny przedział dwumianowy w Julii (Distributions.jl).}]
using Distributions

n, s, alpha = 100, 7, 0.05
phat = s / n
z = quantile(Normal(), 1 - alpha/2)
center = (phat + z^2/(2n))/(1+z^2/n)
half = z/(1+z^2/n)*sqrt(phat*(1-phat)/n + z^2/(4n^2))
wilson = (center-half, center+half)
lower = s == 0 ? 0.0 : quantile(Beta(s, n-s+1), alpha/2)
upper = s == n ? 1.0 : quantile(Beta(s+1, n-s), 1-alpha/2)
println((wilson=wilson, clopper=(lower, upper)))
\end{lstlisting}

\paragraph{Interpretacja kodu.}
Funkcje \texttt{beta.ppf} oraz \texttt{quantile(Beta(...))}
służą wyłącznie jako wygodny sposób numerycznego
rozwiązania równań definiujących Cloppera--Pearsona;
nie oznacza to, że z góry przyjęliśmy bayesowski rozkład
a priori. Kompletne uzasadnienie pokrycia zostało podane
w twierdzeniu~\ref{thm:clopper-pearson}.


% ================================================================
% ROZDZIAL 12 -- WIELOKROTNE TESTOWANIE
% ================================================================
